\chapter{Comprehensive Final Review Checklist}\label{ch:final-review-checklist}
This appendix is a compact review map for the full course.
It is not a substitute for the chapter practice questions; it is a last-pass checklist of the concepts that repeatedly connect exam problems, homework, and projects.


\section{Core Topic Checklist}\label{sec:app-e-core-checklist}

\subsection*{Architecture and Performance}
\begin{itemize}
    \item Know the grain sizes of parallelism: bit-level, instruction-level, thread-level, and inter-program parallelism.
    \item Distinguish SIMD, MIMD, packed-word SIMD, and pipelined vector execution.
    \item Use the CPU-time equation $\text{time}=IC\times CPI\times CT$ and know which term an optimization changes.
    \item Apply Amdahl's law for fixed-size speedup and Gustafson's law for scaled workloads.
    \item Use arithmetic means for times and harmonic means for rates.
\end{itemize}

\subsection*{Caches and Memory Hierarchy}
\begin{itemize}
    \item Split addresses into tag, index, and offset fields from $(C,B,S)$ cache notation.
    \item Classify misses as compulsory, capacity, conflict, or coherence misses when multiprocessors are involved.
    \item Connect each cache optimization to the term it improves: hit time, miss ratio, or miss penalty.
    \item Compare write-back/write-through and write-allocate/write-no-allocate behavior.
    \item Explain why VIPT caches must keep index plus offset bits within the page offset unless extra synonym-handling machinery exists.
    \item Treat tiling as a working-set transformation: choose tiles so the active data footprint fits in cache and does not create severe set conflicts.
\end{itemize}

\subsection*{Pipelines and Branch Prediction}
\begin{itemize}
    \item Interpret IF/ID/EX/MEM/WB timing diagrams cycle by cycle.
    \item Identify structural, data, and control hazards and connect them to inserted stalls or squashes.
    \item Explain why forwarding fixes many RAW hazards but does not eliminate the classic load-use stall.
    \item Compare delay slots, squash slots, branch target buffers, return-address stacks, and direction predictors.
    \item Work through one-bit, two-bit, GSelect, GShare, Yeh/Patt, perceptron, and hybrid predictor lookup/update behavior.
\end{itemize}

\subsection*{Dynamic and Compiler-Scheduled ILP}
\begin{itemize}
    \item Distinguish FICI, FICO, and FOCO execution models.
    \item Explain how Tomasulo-style scheduling uses reservation stations, tags, result broadcast, and register renaming.
    \item Explain why precise exceptions require in-order architectural state update, often through a reorder buffer or checkpoint/repair.
    \item Draw data dependence graphs and classify pure, anti-, and output dependences.
    \item Use dependence height as a list-scheduling priority.
    \item Explain VLIW as a trade: simpler dynamic hardware in exchange for stronger compiler scheduling responsibility.
    \item For software pipelining, compute $ResII$, $RecII$, and $MII=\max(ResII,RecII)$, then reason about prologue, kernel, and epilogue code.
\end{itemize}

\subsection*{Multiprocessors, Synchronization, and Consistency}
\begin{itemize}
    \item Distinguish shared-memory and message-passing systems.
    \item Explain coherence as ordering and propagation for one memory location, not as full synchronization correctness.
    \item Trace MSI/MESI/MOESI-style state transitions and know what \texttt{E}, \texttt{O}, and \texttt{F} add.
    \item Distinguish abort/retry from intervene behavior when dirty data is requested by another cache.
    \item Explain false sharing as unnecessary coherence traffic caused by unrelated variables sharing one cache block.
    \item Compare snooping and directory-based coherence, including why limited-pointer directories trade space against overflow behavior.
    \item Explain why test-and-test-and-set reduces coherence traffic compared with repeated test-and-set.
    \item Use locks for mutual exclusion around critical sections and barriers for regrouping a set of threads at a phase boundary.
    \item Define a data race and explain why coherence alone does not remove data races.
    \item Compare strict, sequential, and weak/relaxed consistency models.
    \item Know what store fences, load fences, full fences, acquire operations, and release operations prevent.
\end{itemize}

\subsection*{Interconnection Networks}
\begin{itemize}
    \item Distinguish direct topologies, such as meshes and tori, from indirect topologies, such as crossbars and multi-stage networks.
    \item For $k$-ary $n$-dimensional networks, compute node count as $k^n$.
    \item Define bisectional bandwidth as the aggregate bandwidth crossing the minimum equal-halves cut.
    \item Explain why a torus has more bisectional bandwidth than the corresponding mesh.
    \item Compare deterministic, adaptive, and oblivious routing.
    \item Explain Valiant's algorithm as randomized load balancing through an intermediate node.
    \item Distinguish packet-level and flit-level flow control.
    \item Compare store-and-forward, cut-through, wormhole, and virtual-channel flow control.
\end{itemize}


\section{Review Clarifications}\label{sec:app-e-clarifications}

\subsection*{Tiling}
Tiling is useful when the program has reuse but the untiled working set is too large for the cache.
The transformation changes loop order so the program operates on cache-sized pieces before moving on.
For exam-style problems, the first check is usually the footprint inequality: active arrays, tile dimensions, and bytes per element must fit within the effective cache capacity.
The second check is placement: a tile that fits by byte count can still perform poorly if the active blocks conflict in the same sets.

\subsection*{Iterative Modulo Scheduling}
Modulo scheduling is easiest to understand as loop unrolling followed by rerolling.
Unrolling exposes overlap among different iterations, but unrolling too far makes code large.
Software pipelining keeps the repeating overlapped pattern as a compact kernel and adds a prologue and epilogue to handle startup and drain.

Iterative modulo scheduling begins at the lower bound $MII=\max(ResII,RecII)$.
$ResII$ comes from resource pressure.
$RecII$ comes from loop-carried dependence cycles.
The compiler tries to schedule the loop at that initiation interval; if it cannot, it increments $II$ and tries again.
An operation that would appear in cycle $x$ of the unrolled schedule maps to row $x\bmod II$ of the modulo reservation table.
Operations that conceptually wrap around the compact kernel explain why prologue and epilogue code are needed.

\subsection*{Memory Consistency}
The core consistency trap is that coherence is per-location, while synchronization usually depends on ordering between locations.
In the producer/consumer pattern, a coherent machine can still let the consumer observe a flag update before it observes the data update that the flag was meant to publish.
The fix is not merely better coherence; the fix is an ordering rule around synchronization, implemented through acquire/release semantics or explicit fences.

Strict consistency ties the visible memory order to actual time.
Sequential consistency allows a single legal interleaving that preserves each thread's program order, even if that interleaving is not the real-time order.
Weak or relaxed models allow ordinary accesses to move more freely but require synchronization accesses to carry stronger ordering guarantees.

\subsection*{Barriers and Locks}
A lock protects a critical section: only one thread should execute that protected code at a time.
A barrier is a phase boundary: every participating thread waits until all have arrived, and then all may proceed.
Barriers are natural in bulk-synchronous programs such as GPU-style kernel phases.
They are expensive when overused because fast threads idle while waiting for slow ones.

\subsection*{Bisectional Bandwidth}
Bisectional bandwidth is a topology figure of merit.
For a legal equal-halves cut, sum the bandwidths of the links crossing the cut.
The bisectional bandwidth is the minimum such aggregate over the allowed equal splits.
For regular meshes and tori, the intended cuts are along topology dimensions, not arbitrary diagonal cuts chosen to inflate the count.
This value is a potential maximum, not a guarantee: routing, buffering, and traffic conflicts can keep real throughput below it.


\section{Fast Self-Check Prompts}\label{sec:app-e-self-check}
\begin{enumerate}
    \item Given a cache geometry, can you compute tag/index/offset bits and explain whether a miss is compulsory, capacity, conflict, or coherence?
    \item Given a branch predictor table and branch history, can you show both the prediction and the update?
    \item Given a dependence graph, can you identify false dependences and explain how renaming removes them?
    \item Given a loop body and functional-unit counts, can you compute $ResII$, $RecII$, and $MII$?
    \item Given a coherence trace, can you count memory reads, memory writes, invalidations, and cache-to-cache transfers?
    \item Given a lock implementation, can you explain which part provides atomicity and which part provides ordering?
    \item Given two threads accessing two variables, can you decide which outcomes are legal under strict consistency, sequential consistency, and relaxed ordering?
    \item Given a network topology, can you compute bisectional bandwidth and explain why routing may still perform worse than that theoretical limit?
    \item Given a packet/flit scenario, can you explain whether store-and-forward, cut-through, wormhole, or virtual channels best describes the behavior?
    \item Given the advanced-topics appendix, can you state why quantum machines and superconducting processors are architectural accelerators rather than simple CMOS replacements?
\end{enumerate}
